\documentclass[10pt,a4paper]{amsart}
\usepackage[margin=19mm]{geometry}
\usepackage{amsmath,amssymb,mathtools}
\usepackage[scr=rsfs,cal=euler]{mathalfa}
\usepackage{xfrac}
\usepackage[unicode,hidelinks,bookmarks=false]{hyperref}
\newcommand{\ie}{i.e.\ }
\newcommand{\cf}{cf.\ }
\newcommand{\resp}{respectively\ }
\newcommand{\Subsection}{\subsection}
\providecommand{\nobreakdash}{\nobreak\hbox{-}}
\setlength{\emergencystretch}{2em}
\multlinegap0pt
\usepackage{bm}
\usepackage{bbm}
\usepackage{dsfont}
\usepackage{xspace}
\usepackage{cancel}
\usepackage{mathtools}
\usepackage{amsthm}
\makeatletter
\@ifundefined{theorem}{\newtheorem{theorem}{Theorem}}{}
\@ifundefined{cor}{\newtheorem{cor}[theorem]{Corollary}}{}
\@ifundefined{lemma}{\newtheorem{lemma}[theorem]{Lemma}}{}
\makeatother
\title[Alternating sign matrices with reflective symmetry and plane]{Alternating sign matrices with reflective symmetry and plane partitions: $n+3$ pairs of equivalent statistics and a Cauchy-type identity}
\author{Ilse Fischer, Hans H\"ongesberg}
\date{Source: arXiv:2207.04469v2 (CC BY 4.0)}
\begin{document}
\maketitle

\noindent\emph{Source and licence.} Alternating sign matrices with reflective symmetry and plane partitions: $n+3$ pairs of equivalent statistics and a Cauchy-type identity. Version arXiv:2207.04469v2 (CC BY 4.0), distributed under the Creative Commons Attribution 4.0 International licence, as recorded in the arXiv licence metadata for this exact version and shown on the abstract page https://arxiv.org/abs/2207.04469v2. Full rights notice: licenses/arxiv-2207.04469-CC-BY-4.0.txt.\par\smallskip

\noindent\textit{Editorial scope.} Counted unit: the Cor whose source label is \texttt{jacobitrudi3cor} in the article (source0.tex lines 3568--3597, source label \texttt{jacobitrudi3cor}), delivered together with the article's own complete original proof of it (source0.tex lines 3599--3649, one \texttt{proof} environment of 51 lines). No mathematics is written, repaired or re-derived: statement and proof are the article's own text line for line. Required context retained uncounted: jacobitrudi3lem (lines 3302--3333); lem:jacobitrudi3lem1 (lines 3305--3332); lem:jacobitrudi3lem2 (lines 3307--3317); prefactor (lines 3541--3548). Every definition, display and label of those lines is retained unchanged. Omitted from this package and not redistributed: 1--3301, 3318--3540, 3549--3567, 3598--3598, 3650--3791. Nothing inside a retained excerpt is omitted or altered: every retained body is delivered exactly as the article's lines are written, so no omission receipt of any kind accompanies this package. Citations: the retained excerpts carry no citation command at all, so no bibliography entry of the article is transcribed and no reference is redistributed. Reference adaptation: none was needed and none was made; every reference inside the delivered unit resolves inside it, so no unresolved reference or citation is produced. Printed slips kept literal: no printed word, symbol, hypothesis or number of the counted statement or of its proof was altered. Layout and class: the article's own class is \texttt{amsart} with packages including inputenc, amssymb, url, mathdots, verbatim, mathtools, nccmath, xcolor, geometry, thmtools, caption, tikz, contour, ytableau; the wrapper is the lane's pinned \texttt{amsart} editorial wrapper at 10pt on A4, which restates the theorem-like environments the retained text needs and adds the article's own macro definitions that the retained text needs (none), with extra wrapper packages (bm, bbm, dsfont, xspace, cancel, mathtools). The delivered numbering is the unit's own. Assets: no retained span contains a figure environment, a graphics-inclusion command, a PDF-page inclusion, a table or a TikZ picture; no image file is copied into this package, the package declares no asset dependency and no image file is redistributed with it. Licence: version arXiv:2207.04469v2 of this article is distributed under the Creative Commons Attribution 4.0 International licence, as recorded in the arXiv licence metadata for this exact version and shown on the abstract page https://arxiv.org/abs/2207.04469v2. A full rights notice is delivered at licenses/arxiv-2207.04469-CC-BY-4.0.txt. Characters: every retained excerpt is pure ASCII, so the delivered engine, XeLaTeX with its default Latin Modern text font, reports no missing character.
\begin{lemma} 
	\label{jacobitrudi3lem}
	In the following identities, the argument of all complete homogeneous symmetric functions~$h$ is $(u X_1,\ldots,u X_n,v X_1^{-1},\ldots,v X_n^{-1})$:
	\begin{enumerate}
		\item\label{lem:jacobitrudi3lem1} For $n \ge 1$, the following identity holds: 
		\begin{multline*} 
				\det\limits_{1 \le i,j \le n} \left( \left(u^2 X_i^2 + u w X_i\right)^j + \left(v^2 X_i^{-2} + v w X_i^{-1}\right)^j \right)
				\det\limits_{1 \le i,j \le n} \left( \left(u^2 X_i^2 + u w X_i\right)^j - \left(v^2 X_i^{-2} + v w X_i^{-1}\right)^j \right)\\
			\times \prod_{1 \le i < j \le n} (X_j-X_i)^{-1} (u^{-1} v X_j^{-1}-u^{-1}v X_i^{-1})^{-1} 
				\prod_{i,j=1}^{n}  (u^{-1} v X_j^{-1}-X_i)^{-1} \\
			= \frac{(-1)^n u^{n^2+\binom{n}{2}} v^{-\binom{n}{2}}}{2} 
			\det_{1 \le i, j \le n } \left( \sum_{k=j}^{2j} \binom{j}{k-j} w^{2j-k}  \left( h_{k-i+1} - u^{-i+1+n} v^{-i+1+n} h_{k+i-2n-1} \right) \right) 
			\\
			\times \det_{1 \le i,j \le n} \left(  \sum_{k=j}^{2j} \binom{j}{k-j} w^{2j-k} 
			(h_{k+i-n-1} + u^{i-1} v^{i-1} h_{k-i-n+1} )  \right).
		\end{multline*} 
		\item\label{lem:jacobitrudi3lem2} For $n \ge 1$, the following identity holds:
		\begin{multline*}
			\det\limits_{1 \le i,j \le n} \left( \left(u^2 X_i^2 + u w X_i\right)^{j-1} + \left(v^2 X_i^{-2} + v w X_i^{-1}\right)^{j-1} \right) \\
			\times\det\limits_{1 \le i,j \le n} \left( \left(u^2 X_i^2 + u w X_i\right)^j - \left(v^2 X_i^{-2} + v w X_i^{-1}\right)^j \right) \\
			\times\prod_{1 \le i < j \le n} (X_j-X_i)^{-1} (u^{-1} v X_j^{-1}-u^{-1}v X_i^{-1})^{-1} 
			\prod_{i,j=1}^{n}  (u^{-1} v X_j^{-1}-X_i)^{-1}\\
			= (-1)^n u^{n^2+\binom{n}{2}} v^{-\binom{n}{2}} 
			\det_{1 \le i, j \le n-1 } \left( \sum_{k=j}^{2j} \binom{j}{k-j} w^{2j-k}  \left( h_{k-i} - u^{-i+n} v^{-i+n} 
			h_{k+i-2n} \right) \right) 
			\\
			\times 
			\det_{1 \le i,j \le n} \left(  \sum_{k=j}^{2j} \binom{j}{k-j} w^{2j-k} 
			(h_{k+i-n-1}+ u^{i-1} v^{i-1}  h_{k-i-n+1} \right).
		\end{multline*}
	\end{enumerate} 
\end{lemma} 

	\begin{lemma} 
		\label{prefactor} 
		Let $n \ge 1$ and $0 \le m \le n-1$. Then we have 
		\begin{equation*}
		\sum_{l=1}^n  \frac{(-u X_l - v X_l^{-1})^{m}}{\prod_{j \not= l} (X_j-X_l)(u-v X_l^{-1} X_j^{-1})} =  
		\begin{cases}1, & m=n-1,\\0, & m<n-1. \end{cases}
		\end{equation*} 
	\end{lemma} 

	\begin{cor} 
		\label{jacobitrudi3cor}
		In the following identities, the argument of all complete homogeneous symmetric functions~$h$ is $(u X_1,\ldots,u X_n,v X_1^{-1},\ldots,v X_n^{-1}).$
		\begin{enumerate}
			\item\label{cor:jacobitrudi3cor1} For $n \ge 1$, the following identity holds: 
			\begin{multline*}
				\frac{
					\det_{1 \le i,j \le n} \left( \left(u^2 X_i^2 + u w X_i\right)^j - \left(v^2 X_i^{-2} + v w X_i^{-1}\right)^j \right)}
				{\prod_{1 \le i < j \le n} (X_j-X_i)(v X_i^{-1} X_j^{-1}-u) \prod_{i=1}^{n} (u X_i -v X_i^{-1})} \\
				=
				\frac{1}{2} \det_{1 \le i,j \le n} \left(  \sum_{k=j}^{2j} \binom{j}{k-j} w^{2j-k} 
				(h_{k+i-n-1} + u^{i-1} v^{i-1} h_{k-i-n+1} )  \right). 
			\end{multline*} 
			\item\label{cor:jacobitrudi3cor2}For $n \ge 1$, the following identity holds:
			\begin{multline*}
				\frac{
					\det_{1 \le i,j \le n} \left( \left(u^2 X_i^2 + u w X_i\right)^j + \left(v^2 X_i^{-2} + v w X_i^{-1}\right)^j \right)}
				{\prod_{1 \le i < j \le n} (X_j-X_i)(u- v X_i^{-1} X_j^{-1})} \\ 
				=
				\det_{1 \le i, j \le n } \left( \sum_{k=j}^{2j} \binom{j}{k-j} w^{2j-k}  \left(  h_{k-i+1} - u^{-i+1+n} v^{-i+1+n} h_{k+i-2n-1} \right) \right).
			\end{multline*} 
			\item\label{cor:jacobitrudi3cor3} For $n \ge 1$, the following identity holds:
			\begin{multline*}
				\frac{1}{2}  \frac{\det_{1 \le i,j \le n} \left( \left(u^2 X_i^2 + u w X_i\right)^{j-1} + \left(v^2 X_i^{-2} + v w X_i^{-1}\right)^{j-1} \right)}
				{\prod_{1 \le i < j \le n} (X_j-X_i)(u - v X_i^{-1} X_j^{-1})} \\
				=
				\det_{1 \le i, j \le n-1} \left( \sum_{k=j}^{2j} \binom{j}{k-j} w^{2j-k}  \left( h_{k-i} - u^{-i+n} v^{-i+n} h_{k+i-2n} \right) \right).
			\end{multline*} 
		\end{enumerate} 
	\end{cor} 

	\begin{proof} 
		The proof is by induction on $n$. For small $n$, the identities can be shown by direct computations. We assume that 
		all identities are proved for $n-1$. 
		We start by proving the third identity. Now observe that, by expanding the determinant with respect to the first column and the induction 
		hypothesis applied to the case $n-1$ in \eqref{cor:jacobitrudi3cor2}, we obtain
		\begin{multline} 
			\label{step}
			\frac{1}{2}  \frac{\det_{1 \le i,j \le n} \left( \left(u^2 X_i^2 + u w X_i\right)^{j-1} + \left(v^2 X_i^{-2} + v w X_i^{-1}\right)^{j-1} \right)}
			{\prod_{1 \le i < j \le n} (X_j-X_i)(u - v X_i^{-1} X_j^{-1})} \\
			= \sum_{l=1}^n (-1)^{l+1} 
			\frac{\det_{\substack{1 \le i \le  n, i \not=l  \\ 1 \le j \le n-1}}
				\left( \left(u^2 X_i^2 + u w X_i\right)^{j} + \left(v^2 X_i^{-2} + v w X_i^{-1}\right)^{j} \right)}
			{\prod_{1 \le i < j \le n} (X_j-X_i)(u - v X_i^{-1} X_j^{-1})} \\
			= \sum_{l=1}^n \frac{1}{\prod_{j \not= l} (X_j-X_l)(u-v X_l^{-1} X_j^{-1})} 
			\det_{1 \le i, j \le n-1} \left( \sum_{k=j}^{2j} \binom{j}{k-j} w^{2j-k}  \left( h^{(l)}_{k-i+1} - u^{-i+n} v^{-i+n} h^{(l)}_{k+i-2n+1} \right) \right), 
		\end{multline} 
		where the argument of the complete homogeneous symmetric functions in the $l$-th summand is 
		\begin{equation}
			\label{argument}
			(u X_1,\ldots,\widehat{u X_l},\ldots,u X_n,v X_1^{-1},\ldots,\widehat{v X_l^{-1}},\ldots,v X_n^{-1}),
		\end{equation} 
		which is indicated by $h^{(l)}$. We will use 
		\begin{multline*} 
			h_k(X_1,\ldots,\widehat{X_p},\ldots,\widehat{X_q},\ldots,X_n) = h_k(X_1,\ldots,\widehat{X_q},\ldots,X_n) - X_p \, h_{k-1}(X_1,\ldots,\widehat{X_q},\ldots,X_n) \\
			= h_k(X_1,\ldots,X_n) - (X_p+X_q) h_{k-1}(X_1,\ldots,X_n) + X_p X_q h_{k-2}(X_1,\ldots,X_n).
		\end{multline*} 
		Applied to our argument in \eqref{argument}, we see that 
		\begin{equation*}
		h_k^{(l)}= h_k - (u X_l + v X_l^{-1}) h_{k-1} + u v\,  h_{k-2},
		\end{equation*}
		and, therefore, \eqref{step} is further equal to 
		\begin{multline}
			\label{eq:splittedhsum}
			\sum_{l=1}^n \frac{1}{\prod_{j \not= l} (X_j-X_l)(u-v X_l^{-1} X_j^{-1})} 
			\det_{1 \le i, j \le n-1} \left( \sum_{k=j}^{2j} \binom{j}{k-j} w^{2j-k} \right. \\
			\times ( h_{k-i+1} - u^{-i+n+1} v^{-i+n+1} h_{k+i-2n-1} 
			- (u X_l + v X_l^{-1}) (h_{k-i} - u^{-i+n} v^{-i+n} h_{k+i-2n})  \\
			\left. \phantom{\sum \binom{n}{k}} \left. + u v (h_{k-i-1} - u^{-i+n-1} v^{-i+n-1} h_{k+i-2n+1}) \right) \right).
		\end{multline} 
		Setting 
		\begin{equation*}
		a_{i,j} = \sum_{k=j}^{2j} \binom{j}{k-j} w^{2j-k} (h_{k-i} - u^{-i+n} v^{-i+n} h_{k+i-2n}),
		\end{equation*}
		the determinant is equal to
		\begin{equation*}
		\det_{1 \le i,j \le n-1} \left( a_{i-1,j} - (u X_l + v X_l^{-1}) a_{i,j} + u v a_{i+1,j} \right).
		\end{equation*}
		By the linearity in the rows, we can write each of the $n$ determinants in \eqref{eq:splittedhsum} as a sum of $3^{n-1}$ determinants, by choosing in a particular row $i$ either of the three terms, namely $a_{i-1,j}$, $-(u X_l + v X_l^{-1}) a_{i,j}$ or $u v a_{i+1,j}$. We pull out $u X_l + v X_l^{-1}$ whenever we choose the second term, so that the determinant is independent of $l$. Across the different choices of $l$, we combine the determinants where in each row we have made precisely the same choices among the three terms (all these determinants are of course equal since we pulled out the appropriate power of $u X_l + v X_l^{-1}$). If we have chosen $m$ times the second term, then we have a prefactor that is equal to the left-hand side of the identity in Lemma~\ref{prefactor}. By the lemma, the prefactor is only nonzero if $m=n-1$. This establishes \eqref{cor:jacobitrudi3cor3} for $n$.
		
		Now \eqref{cor:jacobitrudi3cor1} follows from \eqref{cor:jacobitrudi3cor3} and Lemma~\ref{jacobitrudi3lem}~\eqref{lem:jacobitrudi3lem2}, whereas \eqref{cor:jacobitrudi3cor2} finally follows from \eqref{cor:jacobitrudi3cor1} and Lemma~\ref{jacobitrudi3lem}~\eqref{lem:jacobitrudi3lem1}.
	\end{proof} 

\end{document}
